% chapters/15/section-04.tex
\section{Waktu dan Kekuatan Bunyi}
\label{sec:waktu-volume}

\subsection{Bunyi Harus Terdengar pada Kejadiannya}

Indikator ketiga pada bab ini menuntut lebih dari sekadar ``bunyi sudah
terpasang''. Bunyi harus terdengar pada saat yang tepat. Ada tiga hal yang
menentukan.

Pertama, \textbf{bunyi dipanggil pada peristiwa, bukan pada keadaan}. Bunyi
lompat dipanggil ketika tombol lompat ditekan, bukan selama tombol itu
tertahan. Memanggilnya berdasarkan keadaan menghasilkan bunyi yang berulang
setiap bingkai, dan itu terdengar seperti kerusakan.

Kedua, \textbf{perhatikan di mana pemeriksaan dilakukan}. Pada Bab~XIII, tombol
lompat dibaca di \code{Update} lalu dilaksanakan di \code{FixedUpdate}, sebab
lompatan berkaitan dengan fisika. Bunyi tidak berkaitan dengan fisika. Apabila
bunyi baru dipanggil di \code{FixedUpdate}, sedangkan bingkai fisika berjalan
lebih jarang daripada bingkai tampilan, bunyi dapat tertunda sampai dua atau
tiga bingkai. Pada permainan yang berjalan 60 bingkai per detik, tiga bingkai
berarti sekitar 50 milidetik, dan pada bunyi yang pendek, penundaan sebesar itu
sudah terasa. Panggillah bunyi pada peristiwanya dibaca, yaitu di
\code{Update}.

Ketiga, \textbf{jangan menunggu animasi selesai}. Bunyi adalah umpan balik
seketika. Pemain menekan tombol, dan ia perlu segera memperoleh tanggapan bahwa
tekanannya diterima. Menunda bunyi sampai animasi lompatan mencapai bingkai
tertentu membuat permainan terasa lamban, walaupun gerakannya sendiri sudah
benar.

\subsection{Tingkat Kekuatan Bunyi}

Bunyi permainan Anda akan terdengar paling tidak memuaskan pada satu hal:
musik latarnya menenggelamkan efek suaranya. Pemain tidak dapat mendengar
lompatannya sendiri, sehingga setiap lompatan terasa tanpa tanggapan.

Jalan keluarnya bukan mengubah berkas bunyinya satu per satu, melainkan menyusun
tingkat kekuatan bunyi. Susunan yang dipakai pada bab ini diperlihatkan pada
Gambar~\ref{fig:hierarki-volume}.

\begin{figure}[htbp]
\centering
% chapters/15/figures/hierarki-volume.tikz
% Susunan tingkat kekuatan bunyi: satu induk, dua cabang, dan nilai yang
% biasanya dipakai pada setiap tingkat.
\begin{tikzpicture}[
  tingkat/.style={
    rectangle, rounded corners=5pt, draw=ubPrimary, thick, fill=ubAlt,
    align=center, text width=26mm, inner sep=4pt, minimum height=11mm,
    font=\footnotesize
  },
  panah/.style={-{Stealth[length=2.2mm]}, thick, draw=ubAccent},
  nilai/.style={font=\scriptsize, ubPrimary},
  catatan/.style={font=\scriptsize\itshape, align=left}
]

\node[tingkat] (master) at (3.5,1.5) {\code{Master}\\{\scriptsize\itshape seluruh bunyi}};
\node[tingkat] (musik)  at (0.6,-0.9) {\code{Musik}\\{\scriptsize\itshape BGM}};
\node[tingkat] (efek)   at (6.4,-0.9) {\code{Efek}\\{\scriptsize\itshape SFX}};

\draw[panah] (master) -- (musik);
\draw[panah] (master) -- (efek);

\node[nilai, right=1pt] at (master.east) {0 dB};
\node[nilai, below=0.5pt] at (musik.south) {$-12$ dB};
\node[nilai, below=0.5pt] at (efek.south) {$-6$ dB};

\node[catatan, anchor=north west, text width=78mm] at (-0.9,-2.2)
  {Angka pada \istilah{AudioMixer} memakai satuan desibel, yang naik secara
   berlipat, bukan secara lurus. Menurunkan 6 dB berarti mengurangi
   kekuatannya menjadi kira-kira separuh. Karena itu, penggeser pada menu
   pengaturan sebaiknya tetap bernilai lurus dari 0 sampai 1, lalu diubah ke
   desibel melalui perhitungan.};
\end{tikzpicture}

\caption{Tingkat kekuatan bunyi pada \istilah{AudioMixer}}
\label{fig:hierarki-volume}
\end{figure}

Pada \istilah{AudioMixer}, kekuatan bunyi dinyatakan dalam desibel, dan satuannya
tidak naik secara lurus. Menurunkan 6 dB berarti mengurangi kekuatannya menjadi
kira-kira separuh, dan menaikkan 6 dB berarti melipatgandakannya. Karena itu,
penggeser pada menu pengaturan sebaiknya tetap bernilai lurus dari nol sampai
satu, sebagaimana yang dipahami pemain, lalu diubah ke desibel sebelum
diberikan kepada \istilah{AudioMixer}:

\begin{kodekeluaran}
// Penggeser bernilai 0 sampai 1; mixer menuntut desibel.
// Nilai 0 tidak dapat diubah, sebab desibelnya menuju tak hingga negatif.
if (nilaiPenggeser > 0.0001f)
    mixer.SetFloat("VolumeMusik", Mathf.Log10(nilaiPenggeser) * 20f);
else
    mixer.SetFloat("VolumeMusik", -80f);   // dianggap senyap
\end{kodekeluaran}

Nilai $-80$ dB dipakai sebagai pengganti senyap, sebab itulah batas terendah
yang masih diterima \istilah{AudioMixer}.

\begin{tipsbox}
Angka pada Gambar~\ref{fig:hierarki-volume} dapat dipakai sebagai titik awal,
bukan sebagai patokan. Yang perlu Anda lakukan adalah mendengarkan level Anda
sendiri, lalu memutuskan apakah setiap bunyi terdengar pada saat yang tepat.
Mata rantai yang paling sering salah adalah efek suara yang terlalu pelan,
bukan yang terlalu kuat.
\end{tipsbox}

\subsection{Menggeser Nada agar Tidak Terdengar Seperti Mesin}

Ketika pemain berlari di atas lantai, bunyi langkah dipanggil berulang-ulang
dengan cepat. Apabila setiap panggilan memainkan berkas yang sama persis,
telinga mengenali pengulangan itu, dan bunyinya terdengar seperti mesin.
Perbaikannya murah: geser nadanya sedikit setiap kali, dengan besar yang
berbeda-beda.

\begin{kodekeluaran}
// Nada digeser antara 0,92 dan 1,08 kali nada aslinya.
sfxSource.pitch = 1.0f + Random.Range(-0.08f, 0.08f);
sfxSource.PlayOneShot(klipLangkah);
\end{kodekeluaran}

Nilai yang dipakai pada Daftar~\ref{lst:ch15-manajer} adalah delapan persen.
Lebih dari itu, bunyi langkah mulai terdengar berubah-ubah nadanya seperti
mainan; kurang dari itu, pengulangannya masih terdengar. Untuk bunyi yang
jarang dipanggil, misalnya bunyi kalah, penggeser nada justru tidak diperlukan.

\subsection{Meredam Musik untuk Sementara}

Ada kalanya bunyi tertentu perlu terdengar jelas, sedangkan musik latar sedang
ramai. Caranya bukan menghentikan musik, melainkan menurunkannya untuk
sementara, lalu mengembalikannya. Cara itu disebut \istilah{ducking}, dan pada
\istilah{AudioMixer} dilakukan dengan mengubah nilai satu tingkat saja:

\begin{kodekeluaran}
// Musik diturunkan ketika pemain terluka, lalu dinaikkan kembali.
musikMixer.SetFloat("VolumeMusik", -20f);
Invoke(nameof(KembalikanMusik), 0.8f);
\end{kodekeluaran}

\subsection{Empat Percobaan yang Perlu Anda Lakukan}

Nilai pada bab ini adalah titik awal. Empat percobaan berikut memperlihatkan
pengaruh setiap pengaturan, dan semuanya cukup dilakukan dengan mengubah satu
nilai lalu menjalankan kembali permainannya.

\begin{table}[htbp]
\centering
\caption{Empat percobaan pengaturan bunyi}
\label{tab:percobaan-bunyi}
{\setlength{\tabcolsep}{4pt}\footnotesize%
\begin{tabularx}{\textwidth}{@{}L{3.6cm} L{3.4cm} Y@{}}
\toprule
\headrow \textbf{Yang diubah} & \textbf{Menjadi} & \textbf{Yang akan Anda dengar} \\
\midrule
\code{musicVolume} & 0,2 lalu 1,0 &
Pada nilai rendah, musik hampir tidak terdengar dan efek suara mendominasi;
pada nilai tinggi, bunyi lompat mulai tertelan musik. \\
\addlinespace
\code{pitchVariation} & 0,0 lalu 0,25 &
Pada nilai nol, bunyi berulang terdengar seperti mesin. Pada nilai terlalu
besar, bunyinya berubah-ubah seperti mainan. \\
\addlinespace
\code{PlayOneShot} diganti \code{Play} &
Pada skrip langkah &
Dua langkah yang berdekatan saling memotong, sehingga bunyinya terputus-putus. \\
\addlinespace
\istilah{Spatial Blend} pada bunyi benda &
0 lalu 1 &
Pada nilai nol, bunyi benda di sudut level terdengar sama kuatnya dengan yang
ada di dekat pemain. \\
\bottomrule
\end{tabularx}}
\end{table}
